\begin{enumerate}
  \item \textbf{Label silver.} Sebanyak \val{ts.origin.publisher_section_jsonld.pct} artikel set uji berlabel rubrik JSON-LD dan \val{ts.origin.publisher_category.pct} berlabel kategori penerbit. Pemetaannya konservatif dan tervalidasi terhadap pohon IPTC, tetapi tidak dianotasi ulang oleh manusia. Pemeriksaan tidak langsung pada Subbagian~\ref{sec:res-validity} menunjukkan bahwa, di bahasa yang sama, kesepakatan \eng{teacher} dengan label silver tidak lebih rendah daripada dengan label manusia. Walau begitu, skor absolut pada label silver sebaiknya dibaca sebagai estimasi yang mengandung derau, terutama untuk label yang batasnya kabur (\eng{society}, \eng{human interest}, \eng{lifestyle}).
  \item \textbf{Label MasakhaNEWS dianotasi manusia dengan skemanya sendiri}, lalu dipetakan ke IPTC. Hanya MN-DS yang dianotasi langsung dengan topik IPTC, dan dataset itu hanya berbahasa Inggris. \eng{Test set} manusia IPTC multibahasa dari Kuzman dan Ljube\v{s}i\'c~\cite{kuzman2025llm} tidak publik, sehingga tidak dipakai.
  \item \textbf{Cakupan label per bahasa timpang.} Median jumlah label per bahasa adalah \val{ts.lang_labels.median}. Bahasa Afrika hanya mencakup empat sampai lima topik besar. Macro-F1 dihitung atas label yang hadir, sehingga skor per bahasa tidak dapat dibandingkan secara langsung antarbahasa dengan cakupan label berbeda.
  \item \textbf{Keragaman penerbit terbatas pada sebagian bahasa.} Sebanyak \val{ts.div.satu.langs} bahasa praktis berasal dari satu penerbit: BBC/VOA untuk bahasa Afrika, L3Cube untuk bahasa India, dan satu situs untuk beberapa bahasa CC-NEWS. Klaim untuk bahasa-bahasa ini terbatas pada penerbit tersebut. \eng{Site-cluster bootstrap} memperlebar interval kepercayaan sesuai kondisi ini.
  \item \textbf{Batas atas \eng{student} adalah \eng{teacher}.} Kesalahan dan bias \eng{teacher} ikut tertiru. Set \eng{fidelity} mengukur kemiripan dengan \eng{teacher}, bukan kebenaran.
  \item \textbf{Satu perangkat \eng{benchmark}.} Angka absolut di server akan berbeda, tetapi rasio antarmodel yang diukur dalam sesi dan protokol yang sama tetap informatif.
  \item \textbf{Keadilan antarbahasa dan dampak.} Akurasi StaticKD jauh lebih rendah pada bahasa Afrika yang tidak terwakili di data distilasi maupun pra-latih (Subbagian~\ref{sec:discussion}). Karena itu, penerapan sebagai penyaring topik dapat membuat berita dalam bahasa-bahasa ini lebih sering salah dikategorikan atau terbuang. Kami menyarankan agar akurasi dipantau per bahasa sebelum model dipakai, dan agar artikel berkeyakinan rendah dikirim ke \eng{teacher} lewat mode \eng{cascade}.
  \item \textbf{Cakupan tugas.} Penelitian ini hanya mencakup label IPTC level atas dan klasifikasi \eng{single-label}. Level yang lebih dalam dan artikel multi-topik belum diuji.
\end{enumerate}
